\section{Invalidation Semantics}
\label{app:semantics}

% Formal statements for Section 6, following implementation.md section 11 (the
% executable contract). Definitions first, then the fixed point, then the
% termination proposition whose proof Section 6.4 defers here.

\begin{definition}[Snapshot]
  A snapshot \snapshot{} is the content identity of the workspace: the git
  tree hash, with submodule and LFS pointers taken as tree content. Snapshots
  are immutable. An action either produces a new snapshot or leaves the
  current one unchanged.
\end{definition}

\begin{definition}[Gate key]
  For a gate \(g\) with resolved specification \(s_g\) and a snapshot
  \snapshot{}, the gate key is \(\gatekey(g, \snapshot) = H(\snapshot, s_g)\)
  for a collision-resistant hash \(H\). A diff-scoped gate additionally
  includes the base revision \(b\): \(\gatekey(g, \snapshot) = H(b, \snapshot,
  s_g)\).
\end{definition}

\begin{definition}[Ledger]
  The ledger is a partial function
  \(\ledger : (g, \gatekey) \rightharpoonup \{\textsc{pass}, \textsc{fail}(e)\}\),
  where \(e\) is evidence. Absence is \textsc{unknown}, and \textsc{unknown}
  is what schedules a gate.
\end{definition}

\begin{definition}[Invalidation]
  Let an action transform \snapshot{} into \(\snapshot'\) with changed path
  set \(P\), and let \(I(g)\) be the glob set declared by gate \(g\), which
  defaults to \(\{\texttt{**}\}\). The result of \(g\) carries forward to
  \(\snapshot'\) if and only if \(P \cap \mathrm{paths}(I(g)) = \emptyset\);
  otherwise it is absent at \(\snapshot'\). If \(P\) is reported truncated,
  \(P\) is taken to be every path. The rule applies to \textsc{fail} results
  exactly as to \textsc{pass}.
\end{definition}

\begin{definition}[Fixed point]
  Let a task have driver \(d\), blocking gates \(B\) and non-blocking gates
  \(N\). The task is at its fixed point at snapshot \snapshot{} if and only if
  \(d\) has completed at least once, every \(g \in B\) has
  \(\ledger(g, \gatekey(g, \snapshot)) = \textsc{pass}\), and every
  \(g \in N\) either has a result at \snapshot{} or has exhausted its
  infrastructure retries.
\end{definition}

\begin{definition}[Oscillation]
  Let \(\snapshot_0, \ldots, \snapshot_k\) be the lineage of distinct
  snapshots a task has occupied, with \(\snapshot_0\) the base. An action
  producing \(\snapshot' \neq \snapshot_k\) oscillates if and only if
  \(\snapshot' = \snapshot_i\) for some \(i < k\). An action producing
  \(\snapshot' = \snapshot_k\) is a no-op and is not oscillation.
\end{definition}

\begin{proposition}[Bounded termination]
  Let \(R\) be the cap on driver runs, \(F\) the number of gates that name a
  fix action, \(A\) the per-gate cap on fix attempts, \(G\) the number of
  gates, and \(E\) the cap on consecutive infrastructure failures of one step.
  Every task reaches its fixed point or escalates after at most
  \[
    \bigl( (R + FA) + G\,(R + FA + 1) \bigr)\,(E + 1) + c
  \]
  executions, for a small constant \(c\).
\end{proposition}

\begin{proof}[Proof sketch]
  Each decision either terminates (\textsc{done}, \textsc{escalate}) or
  schedules one execution that strictly consumes a bounded quantity. Driver
  completions are bounded by \(R\) and fix completions by \(FA\), so completed
  actions number at most \(R + FA\). Each new snapshot requires a completed
  action, so the lineage has at most \(R + FA + 1\) entries. A gate re-runs
  only after a new snapshot has cleared its result, so gate completions are
  bounded by \(G\,(R + FA + 1)\). An infrastructure failure increments a
  per-step counter that only a success of that step resets, so each step
  suffers at most \(E\) failures per success and \(E\) more before escalating,
  which gives the factor \(E + 1\). Oscillation, the wall-clock budget and the
  token budget only tighten the bound. The implementation's property suite
  asserts this arithmetic directly against generated execution sequences.
\end{proof}

\todo{Tighten the constant \(c\) once the property tests pin it.}
